\documentclass[12pt,a4paper]{letter}
\usepackage[utf8]{inputenc}
\usepackage[T1]{fontenc}
\usepackage{lmodern}
\usepackage[margin=1in]{geometry}

\signature{Ian Todd\\Sydney Medical School\\University of Sydney}
\address{Ian Todd\\Sydney Medical School\\University of Sydney\\Sydney, NSW, Australia\\itod2305@uni.sydney.edu.au}

\begin{document}

\begin{letter}{Editor\\Foundations of Physics}

\opening{Dear Editor,}

I am writing to submit the manuscript ``Decoherence as Dimensional Collapse: The Geometry of the Quantum-Classical Transition'' for consideration in \textit{Foundations of Physics}.

We prove that the dephasing channel maps the full quantum state manifold onto the classical probability simplex, annihilating the $d^2 - d$ off-diagonal (coherence) directions while preserving the $d - 1$ population dimensions. At full-rank diagonal states, the differential of the dephasing map is the Bures-orthogonal projector onto the classical tangent subspace: classical and quantum degrees of freedom are metrically perpendicular. As our main new result, we compute the exact Bures principal angle between the diagonal and off-diagonal tangent subspaces at full-rank qubit states $(|\mathbf{r}| < 1)$, obtaining a closed-form formula that reveals the geometry of the classical--quantum split across the Bloch ball interior. The split is oblique exactly when the state simultaneously carries nonzero coherence and nonzero population imbalance; dephasing monotonically restores orthogonality, so the metric separation between classical and quantum information is state-dependent and sharpened by the decoherence process. The separation angle carries estimation-theoretic meaning: it determines the fraction of SLD Fisher information for the population parameter that survives elimination of coherence nuisance parameters. We generalize to $d$-dimensional systems with two exact results: a spectator theorem proving that for all $d \geq 3$, Bures-orthogonality fails whenever a nonzero coherence pair shares a common decoupled index, and a perturbative cross-Gram formula quantifying the departure from orthogonality near the classical simplex. We also derive the free energy decomposition showing that, when the Hamiltonian is diagonal in the dephasing basis, coherences store non-equilibrium free energy $k_BT \cdot C_r$, which decoherence irreversibly destroys.

The paper connects Zurek's decoherence program---pointer states, einselection, and quantum Darwinism---to the mathematical framework of information geometry in a way that, to our knowledge, has not been done before. The results use standard quantum mechanics alone; no new assumptions, interpretive commitments, or speculative physics are required.

I believe this work is well suited to \textit{Foundations of Physics} because it addresses foundational questions about the quantum-to-classical transition using rigorous mathematical tools from quantum information geometry.

The manuscript has not been submitted elsewhere and is not under consideration at any other journal. There are no conflicts of interest.

\closing{Sincerely,}

\end{letter}
\end{document}
